\documentclass{article}
\usepackage{amsmath, amssymb}

\begin{document}

\title{Pushing--Medium Electromagnetic Field Equations (v1)}
\author{}
\date{}
\maketitle

%%%%%%%%%%%%%%%%%%%%%%%%%%%%%%%%%%%%%%%%%%%%%%%%%%%%%%%%%%%%%%%%%%%%%%%%%%%%%%%
\section*{1. Overview}
%%%%%%%%%%%%%%%%%%%%%%%%%%%%%%%%%%%%%%%%%%%%%%%%%%%%%%%%%%%%%%%%%%%%%%%%%%%%%%%

Electromagnetism in the Pushing--Medium (PM) framework arises from additional
excitations of the same compressible, advective medium that generates the
gravitational fields $\phi$ and $\mathbf u$.  The PM medium supports two
electromagnetic (EM) modes:


\[
\psi(\mathbf r,t) \quad \text{(scalar EM excitation)},
\qquad
\mathbf w(\mathbf r,t) \quad \text{(vector EM excitation)}.
\]


The scalar field $\psi$ represents compressive electric-like distortions, while
the vector field $\mathbf w$ represents vortical or shear magnetic-like
distortions.  These fields propagate through and are advected by the same medium
flow $\mathbf u$ that appears in the gravitational sector
(see \texttt{v1-unified-formulations.tex}).

%%%%%%%%%%%%%%%%%%%%%%%%%%%%%%%%%%%%%%%%%%%%%%%%%%%%%%%%%%%%%%%%%%%%%%%%%%%%%%%
\section*{2. Electromagnetic Sources}
%%%%%%%%%%%%%%%%%%%%%%%%%%%%%%%%%%%%%%%%%%%%%%%%%%%%%%%%%%%%%%%%%%%%%%%%%%%%%%%

Charge density and electric current act as sources for the EM excitations:


\[
S_\psi = 4\pi K_E\,\rho_q,
\qquad
\mathbf S_w = B_q\,\mathbf J_q.
\]


Here $K_E$ and $B_q$ are coupling constants determining the strength of electric
and magnetic interactions in the PM medium.  Observational constraints on these
constants are given in \texttt{electromagnetic-coupling-constraints.tex}.

%%%%%%%%%%%%%%%%%%%%%%%%%%%%%%%%%%%%%%%%%%%%%%%%%%%%%%%%%%%%%%%%%%%%%%%%%%%%%%%
\section*{3. Dynamic PM-Fluid Equations for EM Fields}
%%%%%%%%%%%%%%%%%%%%%%%%%%%%%%%%%%%%%%%%%%%%%%%%%%%%%%%%%%%%%%%%%%%%%%%%%%%%%%%

\subsection*{3.1 Scalar EM Equation (Electric-like)}


\[
\frac{\partial \psi}{\partial t}
+ \nabla\cdot(\psi\,\mathbf u)
=
S_\psi
+
\frac{1}{c_\psi^2}\frac{\partial^2 \psi}{\partial t^2}.
\]


This equation includes:
\begin{itemize}
    \item advection by the medium flow $\mathbf u$,
    \item compressibility via $\nabla\cdot(\psi\mathbf u)$,
    \item finite propagation speed $c_\psi$,
    \item sourcing by charge density.
\end{itemize}

In the static limit ($\partial_t \psi = 0$, $\mathbf u = 0$):


\[
\nabla^2 \psi = -4\pi K_E\,\rho_q,
\]


which is the PM analogue of Coulomb's law.

\subsection*{3.2 Vector EM Equation (Magnetic-like)}


\[
\frac{\partial \mathbf w}{\partial t}
+ (\mathbf u\cdot\nabla)\mathbf w
=
-\nabla \psi
+ \mathbf S_w
+
\frac{1}{c_w^2}\frac{\partial^2 \mathbf w}{\partial t^2}.
\]


This equation includes:
\begin{itemize}
    \item advection by $\mathbf u$,
    \item coupling to the scalar EM field via $-\nabla\psi$,
    \item finite propagation speed $c_w$,
    \item sourcing by electric current.
\end{itemize}

In the static limit:


\[
\nabla^2 \mathbf w = -B_q\,\mathbf J_q,
\]


the PM analogue of Amp\`ere's law.

%%%%%%%%%%%%%%%%%%%%%%%%%%%%%%%%%%%%%%%%%%%%%%%%%%%%%%%%%%%%%%%%%%%%%%%%%%%%%%%
\section*{4. Linear, No-Flow Limit and EM Wave Equations}
%%%%%%%%%%%%%%%%%%%%%%%%%%%%%%%%%%%%%%%%%%%%%%%%%%%%%%%%%%%%%%%%%%%%%%%%%%%%%%%

In the weak-field, no-flow limit ($\mathbf u = 0$) and for small perturbations,
the EM equations reduce to:


\[
\frac{1}{c^2}\frac{\partial^2 \psi}{\partial t^2}
- \nabla^2 \psi
= -4\pi K_E\,\rho_q,
\]



\[
\frac{1}{c^2}\frac{\partial^2 \mathbf w}{\partial t^2}
- \nabla^2 \mathbf w
= -B_q\,\mathbf J_q.
\]


These are PM analogues of the Lorenz-gauge wave equations for scalar and vector potentials.

%%%%%%%%%%%%%%%%%%%%%%%%%%%%%%%%%%%%%%%%%%%%%%%%%%%%%%%%%%%%%%%%%%%%%%%%%%%%%%%
\section*{5. Effective Electromagnetic Fields}
%%%%%%%%%%%%%%%%%%%%%%%%%%%%%%%%%%%%%%%%%%%%%%%%%%%%%%%%%%%%%%%%%%%%%%%%%%%%%%%

Define the effective electric and magnetic fields:


\[
\mathbf E_{\mathrm{eff}}
= -\nabla \psi - \frac{\partial \mathbf w}{\partial t},
\qquad
\mathbf B_{\mathrm{eff}}
= \nabla \times \mathbf w.
\]


These definitions parallel the standard relations between electromagnetic
potentials and fields.

%%%%%%%%%%%%%%%%%%%%%%%%%%%%%%%%%%%%%%%%%%%%%%%%%%%%%%%%%%%%%%%%%%%%%%%%%%%%%%%
\section*{6. Emergence of Maxwell-like Equations}
%%%%%%%%%%%%%%%%%%%%%%%%%%%%%%%%%%%%%%%%%%%%%%%%%%%%%%%%%%%%%%%%%%%%%%%%%%%%%%%

\subsection*{6.1 Faraday's Law}


\[
\nabla \times \mathbf E_{\mathrm{eff}}
= -\frac{\partial \mathbf B_{\mathrm{eff}}}{\partial t}.
\]


\subsection*{6.2 No Magnetic Monopoles}


\[
\nabla\cdot \mathbf B_{\mathrm{eff}} = 0.
\]


\subsection*{6.3 Gauss's Law}

Using the scalar wave equation and a Lorenz-like gauge condition:


\[
\nabla\cdot \mathbf E_{\mathrm{eff}}
= 4\pi K_E\,\rho_q.
\]


\subsection*{6.4 Amp\`ere--Maxwell Law}

Taking the curl of the vector wave equation yields:


\[
\nabla\times \mathbf B_{\mathrm{eff}}
=
\frac{1}{c^2}\frac{\partial \mathbf E_{\mathrm{eff}}}{\partial t}
+ B_q\,\mathbf J_q.
\]


%%%%%%%%%%%%%%%%%%%%%%%%%%%%%%%%%%%%%%%%%%%%%%%%%%%%%%%%%%%%%%%%%%%%%%%%%%%%%%%
\section*{7. Lagrangian for the EM Sector}
%%%%%%%%%%%%%%%%%%%%%%%%%%%%%%%%%%%%%%%%%%%%%%%%%%%%%%%%%%%%%%%%%%%%%%%%%%%%%%%


\[
\mathcal{L}_{\mathrm{EM}}
=
\frac{1}{2c_\psi^2}(\partial_t \psi)^2
- \frac{1}{2}|\nabla \psi|^2
+
\frac{1}{2c_w^2}|\partial_t \mathbf w|^2
- \frac{1}{2}|\nabla \mathbf w|^2
+
\psi\,S_\psi
+
\mathbf w\cdot\mathbf S_w.
\]


The Euler--Lagrange equations reproduce the PM-fluid EM equations in \S3.

%%%%%%%%%%%%%%%%%%%%%%%%%%%%%%%%%%%%%%%%%%%%%%%%%%%%%%%%%%%%%%%%%%%%%%%%%%%%%%%
\section*{8. Coupling to the Gravitational Sector}
%%%%%%%%%%%%%%%%%%%%%%%%%%%%%%%%%%%%%%%%%%%%%%%%%%%%%%%%%%%%%%%%%%%%%%%%%%%%%%%

The EM fields share the same medium flow $\mathbf u$ as the gravitational sector.
In regimes with non-negligible $\mathbf u$ (e.g., near compact objects or in
high-density plasmas), the advection terms $\nabla\cdot(\psi\mathbf u)$ and
$(\mathbf u\cdot\nabla)\mathbf w$ produce mixing between the EM and gravitational
PM fields.  The unified dynamic equations are given in
\texttt{v1-unified-formulations.tex}, \S10.

Observational constraints on the coupling constants $K_E$ and $A_q$ are derived
in \texttt{electromagnetic-coupling-constraints.tex} and
\texttt{coupling-constants-reference.tex}.

%%%%%%%%%%%%%%%%%%%%%%%%%%%%%%%%%%%%%%%%%%%%%%%%%%%%%%%%%%%%%%%%%%%%%%%%%%%%%%%
\section*{9. Summary}
%%%%%%%%%%%%%%%%%%%%%%%%%%%%%%%%%%%%%%%%%%%%%%%%%%%%%%%%%%%%%%%%%%%%%%%%%%%%%%%

The PM electromagnetic sector arises naturally from additional scalar and vector
excitations of the PM medium.  In the linear, no-flow limit, the effective
fields $\mathbf E_{\mathrm{eff}}$ and $\mathbf B_{\mathrm{eff}}$ satisfy
Maxwell-like equations, while the full PM-fluid equations provide a consistent
nonlinear and advective generalization.  Key points:

\begin{itemize}
    \item EM and gravitational fields share the same medium and the same flow $\mathbf u$.
    \item Static limits recover PM analogues of Coulomb and Amp\`ere's laws.
    \item The linear, no-flow limit reproduces the Lorenz-gauge wave equations.
    \item Effective $\mathbf E$ and $\mathbf B$ satisfy all four Maxwell-like equations.
    \item The Lagrangian $\mathcal{L}_{\mathrm{EM}}$ is structurally identical to the
          gravitational Lagrangian with $(\psi,\mathbf w) \to (\phi,\mathbf u)$.
\end{itemize}

\end{document}
